\lecture{Лекция 2}{}

\section{Транспортная задача}

\subsection{Постановка}

Пусть есть $m$ \textbf{пунктов хранения} (поставщики) $A_1, \dots, A_m$ и $n$
\textbf{пунктов назначения} (потребители) $B_1, \dots, B_n$. Заданы:
\begin{itemize}[nosep,left=0pt]
  \item $\{a_i\}_{i=1}^{m}$ — количество единиц груза в каждом пункте хранения
        (запасы);
  \item $\{b_j\}_{j=1}^{n}$ — количество единиц потребности в грузе в каждом
        пункте назначения;
  \item $\norm{c_{ij}}$ — матрица стоимости перевозок из $i$-го в $j$-й пункт
        (тарифы).
\end{itemize}

\medskip
\noindent\textbf{Задача:} составить план перевозок, удовлетворяющий потребности и
минимизирующий общую стоимость.

\medskip
Обозначим через $x_{ij}$ объём перевозок из $i$-го в $j$-й пункт. Тогда
\[
  \begin{cases}
    \displaystyle\sum_{i=1}^{m} \sum_{j=1}^{n} c_{ij}\, x_{ij} \to \min, \\[12pt]
    \displaystyle\sum_{j=1}^{n} x_{ij} \le a_i, \quad \forall i = \overline{1,m}, \\[12pt]
    \displaystyle\sum_{i=1}^{m} x_{ij} = b_j, \quad \forall j = \overline{1,n}, \\[10pt]
    x_{ij} \ge 0, \quad \forall i, j ,
  \end{cases}
\]
где знак «$\le$» в ограничении по запасам стоит у \textbf{открытой} задачи,
а «$=$» — у \textbf{закрытой}.

\inputfig{figures/fig-transport-table}

\medskip
\noindent\textbf{Условие корректности постановки задачи по исходным данным:}
\[
  \text{открытая:} \quad \sum_{i=1}^{m} a_i \ne \sum_{j=1}^{n} b_j ,
  \qquad\qquad
  \text{закрытая:} \quad \sum_{i=1}^{m} a_i = \sum_{j=1}^{n} b_j .
\]

\begin{idea}
Смысл условия: суммарный вывоз не может превысить суммарный запас, а суммарный
завоз обязан покрыть все потребности. Если $\sum a_i = \sum b_j$, то оба
ограничения автоматически обращаются в равенства — весь груз вывезен и все
потребности закрыты. Если $\sum a_i > \sum b_j$, то часть груза остаётся на
складах, и ограничение по запасам должно быть нестрогим. Если же
$\sum a_i < \sum b_j$, то равенства $\sum_i x_{ij} = b_j$ выполнить нельзя
вовсе — допустимое множество пусто, и задачу приходится переформулировать.
\end{idea}

\subsection{Сведение открытой задачи к закрытой}

Для открытой транспортной задачи производится преобразование её к закрытой.

\begin{enumerate}[label=\arabic*),left=0pt]
  \item Пусть $\displaystyle\sum_{i=1}^{m} a_i < \sum_{j=1}^{n} b_j$. Вводим
        \textbf{фиктивного поставщика} $A_{\text{фикт}}$ с запасом
        \[
          a_{\text{фикт}} = \sum_{j=1}^{n} b_j - \sum_{i=1}^{m} a_i ,
          \qquad x_{\text{фикт},\,j} \ge 0 \quad \forall j = \overline{1,n} .
        \]
        Недопоставки будут распределены оптимальным образом.
  \item Аналогично при $\displaystyle\sum_{i=1}^{m} a_i > \sum_{j=1}^{n} b_j$
        вводим \textbf{фиктивного потребителя} $B_{\text{фикт}}$ с потребностью
        \[
          b_{\text{фикт}} = \sum_{i=1}^{m} a_i - \sum_{j=1}^{n} b_j .
        \]
\end{enumerate}

\begin{idea}
Тарифы фиктивной строки (столбца) берут нулевыми: $c_{\text{фикт},\,j} = 0$ —
фиктивная перевозка ничего не стоит и не меняет значения целевой функции.
Поэтому величина $x_{\text{фикт},\,j}$ в оптимальном плане читается как
«сколько единиц груза потребитель $B_j$ недополучил», а
$x_{i,\,\text{фикт}}$ — как «сколько осталось лежать на складе $A_i$».
Размер задачи вырастает на одну строку (столбец), после чего работает обычный
алгоритм для закрытой задачи.
\end{idea}
